\section{Study-level evidence cards}\label{app:cards}
\pending{These cards reflect the 21-study checkpoint; regenerate them after the final study and report-version decisions.}
The item rationales below reproduce the Group A quality register. Bounded outcomes and caveats reproduce the extraction register. They are not a new independent appraisal or a substitute for the cited full texts.
\Needspace{12\baselineskip}
\subsection{GA-S001: Real-time, Work-conserving GPU Scheduling for Concurrent DNN Inference}
\textbf{Preferred report:} ACM Transactions on Computer Systems, 2025; report GA-PX007. \textbf{Current QA:} 8.5/10 (high).
\begin{longtable}{P{0.15\linewidth}P{0.75\linewidth}}
\toprule Extraction field & Recorded observation\\\midrule\endhead
Mechanism & Reset-based temporal preemption + dynamic kernel padding; NVIDIA binary instrumentation\\
Environment & One MI50, 16GB/60CUs; i7-10700, 32GB RAM; Ubuntu20.04, ROCm4.3.0, TVM0.14.0. Separate one GV100, 32GB/80SMs; CUDA11.4, PyTorch1.12.0, TensorRT8.5.1. Driver version NOT EXTRACTED.\\
Workload & DISB A-E, uniform RT except Poisson E; closed-loop BE; ResNet152, VGG19, YOLOv3, BERT-base, GPT2 (128 tokens, GPT2 prefill). RT rates50\%,100\% or10\% solo-throughput according to workload; Apollo/SVL trace. Separate LLM case batch1,100 decode tokens, MLC-LLM.\\
Comparator & RT-only; sequential priority; GPU streams. NVIDIA closed-source strata compare priority-enabled Triton/TensorRT and Orion/PyTorch without dynamic padding.\\
Bounded finding & GV100 PyTorch DISB-C: Orion p99 normalized latency1.17x vs Reef below1.05x; normalized total throughput0.97x vs1.14x(Fig29). TensorRT comparison trades throughput: mean1.21x Triton vs1.11x Reef(Fig28). All ratios retain model-solo normalization; do not mix engines/AMD strata.\\
Transfer limit & Explicit OSDI2022 extension; select journal as preferred report and retain GA-S001 once. Requires deterministic/stateless kernels; closed-source comparison omits padding. Memory assumed sufficient; bandwidth contention and input-shape profiling limits. Printed publication date differs from some issue-year indexing; recorded2025 follows supplied copy and is in cutoff either way.\\
Source anchor & Article pp1:23-1:24, 1:31-1:34; sections6.2, 6.6, 7; Figs28-29; Markdown lines1111-1180, 1600-1820\\
\bottomrule\end{longtable}
\begin{longtable}{P{0.09\linewidth}P{0.10\linewidth}P{0.67\linewidth}}
\toprule QA & Score & Source-based rationale\\\midrule\endhead
Q1 & 1 & RT and BE sharing question explicit\\
Q2 & 1 & Reset/preemption, padding and NVIDIA implementation described\\
Q3 & 0.5 & GPU/CPU/RAM/OS/frameworks clear; driver missing\\
Q4 & 1 & DISB mix, RT relative rates, arrivals and model/token configurations explicit\\
Q5 & 1 & Same-engine closed-source comparisons and priority settings disclosed\\
Q6 & 1 & Latency and throughput denominators explicit; p99 prose anchored toFig29\\
Q7 & 0.5 & Solo reference averaged1000 runs; repeats/duration for all serving tests unresolved\\
Q8 & 0.5 & Captions/numerical comparisons survive; full CDF/plot data absent\\
Q9 & 1 & Section7 details statelessness, determinism, graphs, bandwidth and memory limits\\
Q10 & 1 & Implementation and explicit engine constraints permit interpretation; experiments not rerun\\
\bottomrule\end{longtable}

\Needspace{12\baselineskip}
\subsection{GA-S002: Serving Heterogeneous Machine Learning Models on Multi-GPU Servers with Spatio-Temporal Sharing}
\textbf{Preferred report:} USENIX ATC, 2022; report GA-PS002. \textbf{Current QA:} 8.5/10 (high).
\begin{longtable}{P{0.15\linewidth}P{0.75\linewidth}}
\toprule Extraction field & Recorded observation\\\midrule\endhead
Mechanism & MPS + temporal\\
Environment & RTX2080Ti 11GB; two/server; Ubuntu18.04; CUDA10.2; driver440.64; PyTorch1.10\\
Workload & Nine models; Poisson arrivals; per-model SLO\\
Comparator & Adapted SBP/Nexus; best-fit with/without interference\\
Bounded finding & Local gpulets support interference-aware provisioning; cluster throughput is a different outcome.\\
Transfer limit & Server-level gains cannot be treated as one-GPU utilization.\\
Source anchor & PDF p10 Table2/3; section5; local profiling section3\\
\bottomrule\end{longtable}
\begin{longtable}{P{0.09\linewidth}P{0.10\linewidth}P{0.67\linewidth}}
\toprule QA & Score & Source-based rationale\\\midrule\endhead
Q1 & 1 & Explicit tenant question\\
Q2 & 1 & Mechanism specified\\
Q3 & 1 & Environment completeness\\
Q4 & 1 & Workload specified\\
Q5 & 0.5 & Partial: comparator suitability/configuration\\
Q6 & 1 & Units/latency statistic clear\\
Q7 & 0.5 & Partial: measurement controls\\
Q8 & 1 & Result figure/table\\
Q9 & 0.5 & Partial: validity limits\\
Q10 & 1 & Interpretability; experiments not rerun\\
\bottomrule\end{longtable}

\Needspace{12\baselineskip}
\subsection{GA-S003: PipeSwitch: Fast Pipelined Context Switching for Deep Learning Applications}
\textbf{Preferred report:} USENIX OSDI, 2020; report GA-PS003. \textbf{Current QA:} 9.0/10 (high).
\begin{longtable}{P{0.15\linewidth}P{0.75\linewidth}}
\toprule Extraction field & Recorded observation\\\midrule\endhead
Mechanism & Temporal context switching\\
Environment & V10016GB or T416GB EC2; PyTorch1.3; CUDA10.1\\
Workload & Training batch32; inference8; RN152/InceptionV3/BERT; 100 runs\\
Comparator & Unoptimized/grouped/per-layer transfer\\
Bounded finding & Fig7: inference preemption latency reduced by up to 38.2ms.\\
Transfer limit & Average switch/inference latency is not p99; multi-VM throughput separate.\\
Source anchor & PDF pp11-13; sections6-6.2; Fig7\\
\bottomrule\end{longtable}
\begin{longtable}{P{0.09\linewidth}P{0.10\linewidth}P{0.67\linewidth}}
\toprule QA & Score & Source-based rationale\\\midrule\endhead
Q1 & 1 & Explicit tenant question\\
Q2 & 1 & Mechanism specified\\
Q3 & 0.5 & Partial: environment completeness\\
Q4 & 1 & Workload specified\\
Q5 & 1 & Comparator suitability/configuration\\
Q6 & 1 & Units/latency statistic clear\\
Q7 & 1 & Measurement controls\\
Q8 & 1 & Result figure/table\\
Q9 & 0.5 & Partial: validity limits\\
Q10 & 1 & Interpretability; experiments not rerun\\
\bottomrule\end{longtable}

\Needspace{12\baselineskip}
\subsection{GA-S004: Serving DNNs like Clockwork: Performance Predictability from the Bottom Up}
\textbf{Preferred report:} USENIX OSDI, 2020; report GA-PS004. \textbf{Current QA:} 8.5/10 (high).
\begin{longtable}{P{0.15\linewidth}P{0.75\linewidth}}
\toprule Extraction field & Recorded observation\\\midrule\endhead
Mechanism & Temporal exclusive scheduling\\
Environment & V10032GB; DellR740; worker32cores/768GB; two GPUs/host\\
Workload & Local experiment: one GPU; 15 RN50 copies; 16 clients/model; SLO10-500ms\\
Comparator & Clockwork variants; INFaaS on different hardware\\
Bounded finding & Fig5 measures goodput satisfying SLO, rather than raw request throughput.\\
Transfer limit & Cross-hardware INFaaS comparison is indirect; controller/worker distinction.\\
Source anchor & PDF p10; section6.1; Fig5\\
\bottomrule\end{longtable}
\begin{longtable}{P{0.09\linewidth}P{0.10\linewidth}P{0.67\linewidth}}
\toprule QA & Score & Source-based rationale\\\midrule\endhead
Q1 & 1 & Explicit tenant question\\
Q2 & 1 & Mechanism specified\\
Q3 & 0.5 & Partial: environment completeness\\
Q4 & 1 & Workload specified\\
Q5 & 0.5 & Partial: comparator suitability/configuration\\
Q6 & 1 & Units/latency statistic clear\\
Q7 & 0.5 & Partial: measurement controls\\
Q8 & 1 & Result figure/table\\
Q9 & 1 & Validity limits\\
Q10 & 1 & Interpretability; experiments not rerun\\
\bottomrule\end{longtable}

\Needspace{12\baselineskip}
\subsection{GA-S005: Torpor: GPU-Enabled Serverless Computing for Low-Latency, Resource-Efficient Inference}
\textbf{Preferred report:} USENIX ATC, 2025; report GA-PS005. \textbf{Current QA:} 8.5/10 (high).
\begin{longtable}{P{0.15\linewidth}P{0.75\linewidth}}
\toprule Extraction field & Recorded observation\\\midrule\endhead
Mechanism & API remoting + temporal swapping\\
Environment & Worker48vCPU/384GB; four V10032GB\\
Workload & Eight inference models; warm-up; CV80ms/BERT200ms SLO\\
Comparator & Native; remote GPU; PCIe/NVLink swap\\
Bounded finding & BERT Table4: native42ms, remote43ms, PCIe swap144ms, NVLink swap45ms.\\
Transfer limit & p98 tail; busy time is not SM occupancy; node scaling separate.\\
Source anchor & PDF pp10-11; Table4; sections7.2-7.3\\
\bottomrule\end{longtable}
\begin{longtable}{P{0.09\linewidth}P{0.10\linewidth}P{0.67\linewidth}}
\toprule QA & Score & Source-based rationale\\\midrule\endhead
Q1 & 1 & Explicit tenant question\\
Q2 & 1 & Mechanism specified\\
Q3 & 0.5 & Partial: environment completeness\\
Q4 & 1 & Workload specified\\
Q5 & 1 & Comparator suitability/configuration\\
Q6 & 1 & Units/latency statistic clear\\
Q7 & 0.5 & Partial: measurement controls\\
Q8 & 1 & Result figure/table\\
Q9 & 0.5 & Partial: validity limits\\
Q10 & 1 & Interpretability; experiments not rerun\\
\bottomrule\end{longtable}

\Needspace{12\baselineskip}
\subsection{GA-S006: StreamBox: A Lightweight GPU SandBox for Serverless Inference Workflow}
\textbf{Preferred report:} USENIX ATC, 2024; report GA-PS006. \textbf{Current QA:} 8.0/10 (high).
\begin{longtable}{P{0.15\linewidth}P{0.75\linewidth}}
\toprule Extraction field & Recorded observation\\\midrule\endhead
Mechanism & API/runtime sandbox virtualization\\
Environment & Four V10016GB; 28cores/128GB; CUDA10.1; PyTorch1.3; TensorFlow2.12\\
Workload & Three concurrent inference workflows/GPU\\
Comparator & MPS; INFless; Astraea; Stream-only\\
Bounded finding & Table2 evaluates runtime sandbox workflows on shared GPUs.\\
Transfer limit & Sandbox memory separation is not proof of adversarial security.\\
Source anchor & PDF p10 onward; section9; Table2\\
\bottomrule\end{longtable}
\begin{longtable}{P{0.09\linewidth}P{0.10\linewidth}P{0.67\linewidth}}
\toprule QA & Score & Source-based rationale\\\midrule\endhead
Q1 & 1 & Explicit tenant question\\
Q2 & 1 & Mechanism specified\\
Q3 & 0.5 & Partial: environment completeness\\
Q4 & 1 & Workload specified\\
Q5 & 0.5 & Partial: comparator suitability/configuration\\
Q6 & 1 & Units/latency statistic clear\\
Q7 & 0.5 & Partial: measurement controls\\
Q8 & 1 & Result figure/table\\
Q9 & 0.5 & Partial: validity limits\\
Q10 & 1 & Interpretability; experiments not rerun\\
\bottomrule\end{longtable}

\Needspace{12\baselineskip}
\subsection{GA-S007: KRYPTON: Efficient Performance-Aware GPU Sharing with Compatibility and Isolation through Kernel Space Interception}
\textbf{Preferred report:} USENIX ATC, 2025; report GA-PS007. \textbf{Current QA:} 9.0/10 (high).
\begin{longtable}{P{0.15\linewidth}P{0.75\linewidth}}
\toprule Extraction field & Recorded observation\\\midrule\endhead
Mechanism & Kernel API virtualization + hybrid\\
Environment & A10040GB/RTX4090; Ubuntu20.04; kernel5.15.0-91; driver550.90.07; CUDA12.1; PyTorch2.2.1\\
Workload & Inference/training/rendering in separate containers; solo-throughput targets\\
Comparator & Temporal; MIG; adapted gpulet\\
Bounded finding & Inference-specific sharing comparisons; section2.2 tests fault propagation.\\
Transfer limit & Baseline adaptations and workload mix constrain attribution.\\
Source anchor & PDF p10 Tables1-2; section5; section2.2 Fig3\\
\bottomrule\end{longtable}
\begin{longtable}{P{0.09\linewidth}P{0.10\linewidth}P{0.67\linewidth}}
\toprule QA & Score & Source-based rationale\\\midrule\endhead
Q1 & 1 & Explicit tenant question\\
Q2 & 1 & Mechanism specified\\
Q3 & 1 & Environment completeness\\
Q4 & 1 & Workload specified\\
Q5 & 0.5 & Partial: comparator suitability/configuration\\
Q6 & 1 & Units/latency statistic clear\\
Q7 & 0.5 & Partial: measurement controls\\
Q8 & 1 & Result figure/table\\
Q9 & 1 & Validity limits\\
Q10 & 1 & Interpretability; experiments not rerun\\
\bottomrule\end{longtable}

\Needspace{12\baselineskip}
\subsection{GA-S008: SIRIUS: Colocating ML Inference and Training with Fast GPU Memory Handover}
\textbf{Preferred report:} USENIX ATC, 2025; report GA-PS008. \textbf{Current QA:} 8.5/10 (high).
\begin{longtable}{P{0.15\linewidth}P{0.75\linewidth}}
\toprule Extraction field & Recorded observation\\\midrule\endhead
Mechanism & Temporal + MPS memory handover\\
Environment & V10016GB; Ubuntu20.04; CUDA11.6; PyTorch2.1.2; TVM0.14; LLM case A10080GB\\
Workload & One GPU; six DNNs/up to56 instances + SWIN training; 300s; Poisson MAF; SLO4x solo\\
Comparator & Reimplemented TaskSwitch; SP50/SP75; UM+MPS\\
Bounded finding & Inference p99/SLO outcomes accompany training progress; separate TTFT/TBT LLM case.\\
Transfer limit & TaskSwitch is a reimplementation, not vanilla PipeSwitch.\\
Source anchor & Section6.1 Fig10; section6.6\\
\bottomrule\end{longtable}
\begin{longtable}{P{0.09\linewidth}P{0.10\linewidth}P{0.67\linewidth}}
\toprule QA & Score & Source-based rationale\\\midrule\endhead
Q1 & 1 & Explicit tenant question\\
Q2 & 1 & Mechanism specified\\
Q3 & 0.5 & Partial: environment completeness\\
Q4 & 1 & Workload specified\\
Q5 & 0.5 & Partial: comparator suitability/configuration\\
Q6 & 1 & Units/latency statistic clear\\
Q7 & 1 & Measurement controls\\
Q8 & 1 & Result figure/table\\
Q9 & 0.5 & Partial: validity limits\\
Q10 & 1 & Interpretability; experiments not rerun\\
\bottomrule\end{longtable}

\Needspace{12\baselineskip}
\subsection{GA-S009: Tally: Non-Intrusive Performance Isolation for Concurrent Deep Learning Workloads}
\textbf{Preferred report:} ACM ASPLOS Volume 1, 2025; report GA-PS009. \textbf{Current QA:} 9.0/10 (high).
\begin{longtable}{P{0.15\linewidth}P{0.75\linewidth}}
\toprule Extraction field & Recorded observation\\\midrule\endhead
Mechanism & CUDA API interception + priority\\
Environment & A10040GB; Ubuntu20.04; CUDA12.2; PyTorch2.2; ONNX1.17\\
Workload & Same-GPU pairs; MAF inference traces; high-priority p99\\
Comparator & Time slicing; MPS; priority MPS; TGS\\
Bounded finding & At 10\% idle time: normalized throughput Tally1.19 vs TGS1.65.\\
Transfer limit & Latency protection sacrifices some BE throughput; author manuscript.\\
Source anchor & PDF pp9-11; sections5-5.3; Fig6\\
\bottomrule\end{longtable}
\begin{longtable}{P{0.09\linewidth}P{0.10\linewidth}P{0.67\linewidth}}
\toprule QA & Score & Source-based rationale\\\midrule\endhead
Q1 & 1 & Explicit tenant question\\
Q2 & 1 & Mechanism specified\\
Q3 & 0.5 & Partial: environment completeness\\
Q4 & 1 & Workload specified\\
Q5 & 1 & Comparator suitability/configuration\\
Q6 & 1 & Units/latency statistic clear\\
Q7 & 0.5 & Partial: measurement controls\\
Q8 & 1 & Result figure/table\\
Q9 & 1 & Validity limits\\
Q10 & 1 & Interpretability; experiments not rerun\\
\bottomrule\end{longtable}

\Needspace{12\baselineskip}
\subsection{GA-S010: ParvaGPU: Efficient Spatial GPU Sharing for Large-Scale DNN Inference in Cloud Environments}
\textbf{Preferred report:} IEEE/ACM SC, 2024; report GA-P0589. \textbf{Current QA:} 8.0/10 (high).
\begin{longtable}{P{0.15\linewidth}P{0.75\linewidth}}
\toprule Extraction field & Recorded observation\\\midrule\endhead
Mechanism & MIG + MPS\\
Environment & A10080GB; Ubuntu20.04; CUDA12.0.1; reported PyTorch1.14.0 needs checking\\
Workload & Eleven models; six scenario mixes; same-model MPS inside MIG\\
Comparator & MIG-serving; gpulet and system variants\\
Bounded finding & Fig7 separates internal/external fragmentation; Fig5 counts GPUs.\\
Transfer limit & Author manuscript; unusual runtime version; GPU count is not occupancy.\\
Source anchor & PDF p8 onward; sectionIV; TableIV; Figs5/7\\
\bottomrule\end{longtable}
\begin{longtable}{P{0.09\linewidth}P{0.10\linewidth}P{0.67\linewidth}}
\toprule QA & Score & Source-based rationale\\\midrule\endhead
Q1 & 1 & Explicit tenant question\\
Q2 & 1 & Mechanism specified\\
Q3 & 0.5 & Partial: environment completeness\\
Q4 & 1 & Workload specified\\
Q5 & 0.5 & Partial: comparator suitability/configuration\\
Q6 & 1 & Units/latency statistic clear\\
Q7 & 0.5 & Partial: measurement controls\\
Q8 & 1 & Result figure/table\\
Q9 & 0.5 & Partial: validity limits\\
Q10 & 1 & Interpretability; experiments not rerun\\
\bottomrule\end{longtable}

\Needspace{12\baselineskip}
\subsection{GA-S011: Orion: Interference-aware Fine-grained GPU Sharing for ML Applications}
\textbf{Preferred report:} ACM EuroSys, 2024; report GA-PS011. \textbf{Current QA:} 9.0/10 (high).
\begin{longtable}{P{0.15\linewidth}P{0.75\linewidth}}
\toprule Extraction field & Recorded observation\\\midrule\endhead
Mechanism & CUDA API interception + temporal\\
Environment & V10016GB/CUDA10.2; A10040GB/CUDA11.3; PyTorch1.12; Python3.9\\
Workload & Inference-inference and inference-training; uniform/Poisson; three repeats\\
Comparator & MPS; streams; temporal; REEF-N reimplementation\\
Bounded finding & Interference-aware scheduling compares inference latency and BE throughput.\\
Transfer limit & Prototype clients share a process; REEF-N differs from AMD REEF.\\
Source anchor & PDF pp8-9; sections5.3-6.1; Tables1/3\\
\bottomrule\end{longtable}
\begin{longtable}{P{0.09\linewidth}P{0.10\linewidth}P{0.67\linewidth}}
\toprule QA & Score & Source-based rationale\\\midrule\endhead
Q1 & 1 & Explicit tenant question\\
Q2 & 1 & Mechanism specified\\
Q3 & 0.5 & Partial: environment completeness\\
Q4 & 1 & Workload specified\\
Q5 & 0.5 & Partial: comparator suitability/configuration\\
Q6 & 1 & Units/latency statistic clear\\
Q7 & 1 & Measurement controls\\
Q8 & 1 & Result figure/table\\
Q9 & 1 & Validity limits\\
Q10 & 1 & Interpretability; experiments not rerun\\
\bottomrule\end{longtable}

\Needspace{12\baselineskip}
\subsection{GA-S012: GSLICE: Controlled Spatial Sharing of GPUs for a Scalable Inference Platform}
\textbf{Preferred report:} ACM SoCC, 2020; report GA-PS012. \textbf{Current QA:} 9.0/10 (high).
\begin{longtable}{P{0.15\linewidth}P{0.75\linewidth}}
\toprule Extraction field & Recorded observation\\\midrule\endhead
Mechanism & MPS adaptive provisioning\\
Environment & One V10016GB; XeonGold6148; kernel4.4.0-150; DPDK19.02; CNTK/PyTorch/TensorRT\\
Workload & Two-four inference functions; VGG case SLO50ms; traffic replay\\
Comparator & Standalone functions; MPS provisioning variants\\
Bounded finding & GUE measures efficiency; shadow instance loading takes about12s in Fig15.\\
Transfer limit & GUE is not GPU busy percentage; median latency not p99; academic mirror.\\
Source anchor & PDF pp8-12; section4; Table3; Fig15\\
\bottomrule\end{longtable}
\begin{longtable}{P{0.09\linewidth}P{0.10\linewidth}P{0.67\linewidth}}
\toprule QA & Score & Source-based rationale\\\midrule\endhead
Q1 & 1 & Explicit tenant question\\
Q2 & 1 & Mechanism specified\\
Q3 & 0.5 & Partial: environment completeness\\
Q4 & 1 & Workload specified\\
Q5 & 1 & Comparator suitability/configuration\\
Q6 & 1 & Units/latency statistic clear\\
Q7 & 1 & Measurement controls\\
Q8 & 1 & Result figure/table\\
Q9 & 0.5 & Partial: validity limits\\
Q10 & 1 & Interpretability; experiments not rerun\\
\bottomrule\end{longtable}

\Needspace{12\baselineskip}
\subsection{GA-S013: KRISP: Enabling Kernel-wise RIght-sizing for Spatial Partitioned GPU Inference Servers}
\textbf{Preferred report:} IEEE HPCA, 2023; report GA-P0594. \textbf{Current QA:} 8.0/10 (high).
\begin{longtable}{P{0.15\linewidth}P{0.75\linewidth}}
\toprule Extraction field & Recorded observation\\\midrule\endhead
Mechanism & Spatial partition/API control (AMD)\\
Environment & One MI50/60CUs; dualEPYC7302/512GB; Ubuntu18.04; kernel5.4; ROCm4.5\\
Workload & One/two/four workers and heterogeneous pairs; batch32; aggregate8/16; eight models\\
Comparator & AMD unrestricted concurrency labelled MPSDefault\\
Bounded finding & Kernel-wise right-sizing evaluates throughput and p95 inference latency.\\
Transfer limit & Feature emulated; overhead subtracted; baseline is not NVIDIA MPS.\\
Source anchor & PDF pp9-10; sectionsV-VI; TableIII\\
\bottomrule\end{longtable}
\begin{longtable}{P{0.09\linewidth}P{0.10\linewidth}P{0.67\linewidth}}
\toprule QA & Score & Source-based rationale\\\midrule\endhead
Q1 & 1 & Explicit tenant question\\
Q2 & 1 & Mechanism specified\\
Q3 & 1 & Environment completeness\\
Q4 & 1 & Workload specified\\
Q5 & 0.5 & Partial: comparator suitability/configuration\\
Q6 & 1 & Units/latency statistic clear\\
Q7 & 0.5 & Partial: measurement controls\\
Q8 & 1 & Result figure/table\\
Q9 & 0.5 & Partial: validity limits\\
Q10 & 0.5 & Partial: interpretability; experiments not rerun\\
\bottomrule\end{longtable}

\Needspace{12\baselineskip}
\subsection{GA-S014: XSched: Preemptive Scheduling for Diverse XPUs}
\textbf{Preferred report:} USENIX OSDI, 2025; report GA-PS014. \textbf{Current QA:} 9.0/10 (high).
\begin{longtable}{P{0.15\linewidth}P{0.75\linewidth}}
\toprule Extraction field & Recorded observation\\\midrule\endhead
Mechanism & Temporal API interposition\\
Environment & One GV100 inference case; broader XPU experiments separate\\
Workload & Two BERT-large clients; HP10req/s; saturated BE\\
Comparator & Triton; Triton+priority; Triton+XSched\\
Bounded finding & Fig15: relative p99 1.53/1.51/1.07 respectively versus solo.\\
Transfer limit & Data assumed to fit memory; CUDA graphs and undocumented driver limits.\\
Source anchor & PDF pp14-15; section8 Case3 Fig15; section9\\
\bottomrule\end{longtable}
\begin{longtable}{P{0.09\linewidth}P{0.10\linewidth}P{0.67\linewidth}}
\toprule QA & Score & Source-based rationale\\\midrule\endhead
Q1 & 1 & Explicit tenant question\\
Q2 & 1 & Mechanism specified\\
Q3 & 0.5 & Partial: environment completeness\\
Q4 & 1 & Workload specified\\
Q5 & 1 & Comparator suitability/configuration\\
Q6 & 1 & Units/latency statistic clear\\
Q7 & 0.5 & Partial: measurement controls\\
Q8 & 1 & Result figure/table\\
Q9 & 1 & Validity limits\\
Q10 & 1 & Interpretability; experiments not rerun\\
\bottomrule\end{longtable}

\Needspace{12\baselineskip}
\subsection{GA-S015: Paella: Low-latency Model Serving with Software-defined GPU Scheduling}
\textbf{Preferred report:} ACM SOSP, 2023; report GA-PS015. \textbf{Current QA:} 8.5/10 (high).
\begin{longtable}{P{0.15\linewidth}P{0.75\linewidth}}
\toprule Extraction field & Recorded observation\\\midrule\endhead
Mechanism & Temporal software scheduling + MPS comparison\\
Environment & TeslaT4; XeonSilver4114/64GB; Linux4.15.0; CUDA11.7; modified TVM0.10\\
Workload & Eight models; lognormal sigma1.5/2; stabilized measurements\\
Comparator & CUDA single/multiple streams; MPS; Triton23.03; ablations\\
Bounded finding & p99 latency and throughput under bursty model mixes; Fig12 includes MPS.\\
Transfer limit & Custom Triton TVM wrapper; MPS omitted from Fig11; dynamic batching unsupported.\\
Source anchor & PDF pp9-12; section7 Table3/Figs11-12; section8\\
\bottomrule\end{longtable}
\begin{longtable}{P{0.09\linewidth}P{0.10\linewidth}P{0.67\linewidth}}
\toprule QA & Score & Source-based rationale\\\midrule\endhead
Q1 & 1 & Explicit tenant question\\
Q2 & 1 & Mechanism specified\\
Q3 & 0.5 & Partial: environment completeness\\
Q4 & 1 & Workload specified\\
Q5 & 0.5 & Partial: comparator suitability/configuration\\
Q6 & 1 & Units/latency statistic clear\\
Q7 & 0.5 & Partial: measurement controls\\
Q8 & 1 & Result figure/table\\
Q9 & 1 & Validity limits\\
Q10 & 1 & Interpretability; experiments not rerun\\
\bottomrule\end{longtable}

\Needspace{12\baselineskip}
\subsection{GA-S016: Optimizing High-Throughput Inference on Graph Neural Networks at Shared Computing Facilities with the NVIDIA Triton Inference Server}
\textbf{Preferred report:} Computing and Software for Big Science, 2024; report GA-P0041. \textbf{Current QA:} 8.5/10 (high).
\begin{longtable}{P{0.15\linewidth}P{0.75\linewidth}}
\toprule Extraction field & Recorded observation\\\midrule\endhead
Mechanism & MIG + concurrent model serving\\
Environment & A10080GB; 20/40GB slices; two hosts/four GPUs each; AMD Epyc7543; 100Gbps\\
Workload & Distinct ParticleNet copies within one MIG slice; four requesting workers/model; 120s samples\\
Comparator & 20GB vs40GB slices; scaling settings\\
Bounded finding & Fig10: thrashing after two models/20GB or five/40GB.\\
Transfer limit & Local slice evidence only; multi-server scaling separate; isolation properties described, not attack-tested.\\
Source anchor & HTML Multi-model Scaling/Fig10; Triton implementation; limitations\\
\bottomrule\end{longtable}
\begin{longtable}{P{0.09\linewidth}P{0.10\linewidth}P{0.67\linewidth}}
\toprule QA & Score & Source-based rationale\\\midrule\endhead
Q1 & 1 & Explicit tenant question\\
Q2 & 1 & Mechanism specified\\
Q3 & 0.5 & Partial: environment completeness\\
Q4 & 1 & Workload specified\\
Q5 & 0.5 & Partial: comparator suitability/configuration\\
Q6 & 1 & Units/latency statistic clear\\
Q7 & 0.5 & Partial: measurement controls\\
Q8 & 1 & Result figure/table\\
Q9 & 1 & Validity limits\\
Q10 & 1 & Interpretability; experiments not rerun\\
\bottomrule\end{longtable}

\Needspace{12\baselineskip}
\subsection{GA-S017: VastPipe: A High-Throughput Inference System via Adaptive Space-Division Multiplexing for Diverse Accelerators}
\textbf{Preferred report:} Journal of Computer Science and Technology, 2025; report GA-P0006. \textbf{Current QA:} 6.5/10 (moderate).
\begin{longtable}{P{0.15\linewidth}P{0.75\linewidth}}
\toprule Extraction field & Recorded observation\\\midrule\endhead
Mechanism & MPS adaptive quotas; separate AMD CU-mask stratum\\
Environment & Local NVIDIA A100: TensorRT serving, MPS; AMD MI100: TVM/ROCm with hipExtStreamCreateWithCUMask. GPU memory, OS, driver and runtime versions NOT EXTRACTED for these GPU experiments. NPU Table3 is not the GPU configuration.\\
Workload & Co-located CV pairs (VGG16/GoogLeNet, VGG19/GoogLeNet, ResNet101/InceptionV3, VGG19/ResNet50); NLP/CV pairs with BERT/GPT2; Poisson-generated request mixes. Exact GPU batch sizes and rates NOT EXTRACTED.\\
Comparator & No MPS; default MPS without per-service quotas; RL-selected MPS quotas. AMD no/default CU mask separately.\\
Bounded finding & A100 VGG16+GoogLeNet: throughput 2.7x no MPS and 1.5x default MPS; not a universal average or p99 result. NLP/CV mean 2.43x no MPS, 1.45x default MPS. MI100 mean 1.89x no mask, 1.24x default mask.\\
Transfer limit & Include only local GPU inference evidence, not the main 1000-NPU cluster results. Offline Q-table assumes relatively stable request mix; GPU-specific quotas, repetitions and variance unresolved. Chart bars are not preserved in Markdown. Preliminary ICCD2021 report is linked as an unassessed version lead.\\
Source anchor & Published pp457-459; sections6.7-6.9; Figs14-16; uploaded Markdown lines1150-1290\\
\bottomrule\end{longtable}
\begin{longtable}{P{0.09\linewidth}P{0.10\linewidth}P{0.67\linewidth}}
\toprule QA & Score & Source-based rationale\\\midrule\endhead
Q1 & 1 & Local model co-location question explicit in6.7/6.8\\
Q2 & 1 & Per-model MPS control and CU-mask API specified\\
Q3 & 0.5 & GPU models/backends known; GPU OS/driver/version/memory unresolved\\
Q4 & 1 & Named co-located models and request distribution; exact load incomplete\\
Q5 & 0.5 & No/default/tuned sharing compares same platform; exact quotas incomplete\\
Q6 & 1 & Throughput ratios and requests/s labelled; no tail-latency inference\\
Q7 & 0 & GPU-specific warm-up, run duration, repetitions and variance not recoverable from supplied text\\
Q8 & 0.5 & Figures/captions and numerical prose survive; chart bars do not\\
Q9 & 0.5 & Stable-mix/Q-table assumptions stated; GPU threats incompletely discussed\\
Q10 & 0.5 & GPU results interpretable with caveats; insufficient configuration for replication\\
\bottomrule\end{longtable}

\Needspace{12\baselineskip}
\subsection{GA-S018: ArkGPU: enabling applications' high-goodput co-location execution on multitasking GPUs}
\textbf{Preferred report:} CCF Transactions on High Performance Computing, 2023; report GA-P0037. \textbf{Current QA:} 7.0/10 (moderate).
\begin{longtable}{P{0.15\linewidth}P{0.75\linewidth}}
\toprule Extraction field & Recorded observation\\\midrule\endhead
Mechanism & CUDA API interception + NVML monitoring and admission throttling\\
Environment & Machine with four V100s; Xeon Platinum8163; CentOS7.9.2009; driver460.106.00; CUDA11.1; cuDNN7.6.5. Per-GPU memory/framework versions NOT EXTRACTED.\\
Workload & Five MLPerf inference LC models (ResNet50, SSD-MobileNet, SSD-ResNet34, DLRM, BERT), each paired locally with six BE training models (30 pairs). Single-stream query size1; p99 QoS threshold twice solo. Separate four-GPU app-mixes are not local occupancy evidence.\\
Comparator & MPS (Non-ArkGPU); k8s-native for QPS/utilization; GaiaGPU and adjustment-unit ablation. MPS quota settings NOT EXTRACTED.\\
Bounded finding & Section5.2 reports all30 tested pairs satisfy the 2x-solo p99 QoS target with ArkGPU. Low-contention pairs can lose about5\% QPS from throttling. Four-GPU app-mix goodput gains94.98\% vs k8s-native and38.65\% vs MPS are separate system-level results.\\
Transfer limit & Mixed inference/training qualifies via explicit inference p99/QPS and workload transfer rationale; it does not establish inference-only performance. Polling/API sleeps can add delay; BE checkpoint/restart affects costs. Huge goodput ratios can have near-zero baseline yield. Repetitions/variance and chart coordinates unresolved.\\
Source anchor & Published pp313-319; sections4.3, 5.1-5.4; Table3; Figs14-21; Markdown lines920-1360\\
\bottomrule\end{longtable}
\begin{longtable}{P{0.09\linewidth}P{0.10\linewidth}P{0.67\linewidth}}
\toprule QA & Score & Source-based rationale\\\midrule\endhead
Q1 & 1 & LC protection during BE co-location explicit\\
Q2 & 1 & CUDA interception, NVML polling, sleep and adjustment described in4.3\\
Q3 & 1 & Table3 specifies GPU/CPU/OS/driver/CUDA/cuDNN; memory not transcribed\\
Q4 & 1 & MLPerf single-stream, query size1, 30 pairs and QoS rule explicit\\
Q5 & 0.5 & MPS/native/GaiaGPU compared; quotas and baseline implementation incomplete\\
Q6 & 1 & p99, average, QPS and yield-based goodput distinguished\\
Q7 & 0 & Warm-up/window/repetitions/uncertainty not recoverable in supplied text\\
Q8 & 0.5 & Tables and numerical prose readable; plot data unavailable\\
Q9 & 0.5 & Throttling overhead and low-contention regressions discussed; limited threats analysis\\
Q10 & 0.5 & Algorithm and data-on-request statement; workload/configuration gaps restrict replication\\
\bottomrule\end{longtable}

\Needspace{12\baselineskip}
\subsection{GA-S019: Dyna-P: placement-aware dynamic partitioning for lightweight applications with modern GPUs}
\textbf{Preferred report:} Cluster Computing, 2025; report GA-P0014. \textbf{Current QA:} 7.0/10 (moderate).
\begin{longtable}{P{0.15\linewidth}P{0.75\linewidth}}
\toprule Extraction field & Recorded observation\\\midrule\endhead
Mechanism & MIG dynamic placement/reconfiguration with reconfiguration-last policy\\
Environment & Two A30s, 24GB each; CUDA12.2; driver535.171.04; DCGM3.1.3. Local sections4.2/4.3.1 use one A30. StaCon profiles2g.12gb+1g.6gb+1g.6gb. OS/backend versions NOT EXTRACTED.\\
Workload & 50 SLM inference jobs; batch1-64; DistilGPT, Codegen, Flan-T5-Large, Codellama and Phi-1 as labelled in Table2. Local batch sensitivity:10 inference jobs, batches1-8 vs32-64 on one A30. Later160-A100 cluster study is simulation.\\
Comparator & Full-GPU sequential baseline; fixed StaCon MIG; pre-profiled MISO adapted to A30; ORACLE with future job/size/time knowledge.\\
Bounded finding & Single-A30 batch experiments: smaller partitions increase concurrency but can increase mean JCT; not online p99 evidence. Fig8 broader job set reports39.03\% lower makespan vs StaCon, 12.14\% vs MISO and14.7\% higher normalized STP vs StaCon. Prediction-error sensitivity:15\% perturbation adds8.09\% makespan.\\
Transfer limit & Retain one-GPU inference strata separately from two-GPU aggregate and cluster simulation. JCT/makespan/STP cannot replace request-tail latency or goodput. Reconfiguration can stop/restart active jobs; neither architectural MIG claims nor OOM avoidance prove hostile isolation. Model labels/parameter counts and corrupted overhead fraction require caution.\\
Source anchor & PDF pp11-16; sections4.1-4.3.2; Table2; Figs8-13; Markdown lines700-865\\
\bottomrule\end{longtable}
\begin{longtable}{P{0.09\linewidth}P{0.10\linewidth}P{0.67\linewidth}}
\toprule QA & Score & Source-based rationale\\\midrule\endhead
Q1 & 1 & Fragmentation/concurrent inference question explicit\\
Q2 & 1 & MIG profiles, merge/split and placement algorithms described\\
Q3 & 0.5 & GPU/memory/CUDA/driver/DCGM clear; OS/backend unresolved\\
Q4 & 1 & SLMs, 50 jobs, batch ranges and local10-job strata explicit\\
Q5 & 0.5 & MISO pre-profiled and adapted; ORACLE privileged information disclosed\\
Q6 & 1 & JCT/makespan/normalized STP distinct; avoid p99 interpretation\\
Q7 & 0.5 & Prediction perturbation and power-cap sensitivities; repetitions/variance unresolved\\
Q8 & 0.5 & Table2 and numeric prose readable; plot coordinates missing\\
Q9 & 0.5 & Reconfiguration cost/architecture and prediction limits discussed; incomplete threats treatment\\
Q10 & 0.5 & Implementation described but missing runtime/control detail restricts replication\\
\bottomrule\end{longtable}

\Needspace{12\baselineskip}
\subsection{GA-S020: iGniter: Interference-Aware GPU Resource Provisioning for Predictable DNN Inference in the Cloud}
\textbf{Preferred report:} IEEE Transactions on Parallel and Distributed Systems, 2023; report GA-PX005. \textbf{Current QA:} 8.5/10 (high).
\begin{longtable}{P{0.15\linewidth}P{0.75\linewidth}}
\toprule Extraction field & Recorded observation\\\midrule\endhead
Mechanism & MPS active-thread quotas + interference-aware batch/provisioning and shadow process\\
Environment & Ten EC2 p3.2xlarge instances, one V100 each, 8vCPU/61GB; Triton2.12.0, TensorRT8.0.1.6. MPS active-thread percentage; allocation unit2.5\%. Heterogeneous extension adds20 T4 instances. OS/driver/CUDA version NOT EXTRACTED.\\
Workload & AlexNet, ResNet50, VGG19(ImageNet), SSD(VOC2012); 12 workload/SLO configurations; constant arrivals. Local experiment co-locates four models at25\% quota, batch3. p99 provisioning measured over windows e.g.30s; model-validation experiments repeated3 times with SD bars.\\
Comparator & FFD+; GSLICE+ with iGniter placement; gpu-lets+ modified batch strategy and best-fit placement. These are adapted baselines.\\
Bounded finding & Local four-model latency prediction error1.53\%-5.02\%(Fig13). Fig14 deployment: iGniter6 GPUs/0 SLO violations vs gpu-lets+8/3, FFD+5/10, GSLICE+6/3; 25\% lower hourly instance cost vs gpu-lets+, not25\% higher local utilization. Fig17 shows an initial p99 violation before shadow-process recovery.\\
Transfer limit & Constant arrivals and four vision models limit bursty/LLM transfer. Negligible PCIe contention assumption can underpredict AlexNet; baselines modified. Steady-state/post-correction SLO success is not an unconditional guarantee; shadow standby incurs capacity/memory costs. Chart bars/latency traces unavailable in Markdown.\\
Source anchor & Published pp820-824; sections4.2, 5.1-5.4; Figs13-18; Markdown lines938-1245\\
\bottomrule\end{longtable}
\begin{longtable}{P{0.09\linewidth}P{0.10\linewidth}P{0.67\linewidth}}
\toprule QA & Score & Source-based rationale\\\midrule\endhead
Q1 & 1 & Inference interference and SLO provisioning explicit\\
Q2 & 1 & Active-thread quotas, batching, placement and shadow recovery described\\
Q3 & 0.5 & EC2 GPU/CPU/memory and serving versions clear; OS/driver missing\\
Q4 & 1 & Models, datasets, constant arrivals, quotas and local batch3 clear\\
Q5 & 0.5 & Three adapted baselines and modifications disclosed; effect confounded with placement/batching\\
Q6 & 1 & p99, SLO throughput, cost and prediction error distinguished\\
Q7 & 1 & Model validation three repetitions/SD; p99 windows and update rule explicit; not all experiments have same controls\\
Q8 & 0.5 & Figure captions/numeric prose survive; error bars and full traces not digitizable\\
Q9 & 1 & PCIe assumption, model error and recovery described\\
Q10 & 1 & Public code URL, algorithm and coefficients aid interpretation; code not rerun\\
\bottomrule\end{longtable}

\Needspace{12\baselineskip}
\subsection{GA-S021: InferFair: Towards QoS-aware scheduling for performance isolation guarantee in heterogeneous model serving systems}
\textbf{Preferred report:} Future Generation Computer Systems, 2024; report GA-PX006. \textbf{Current QA:} 7.5/10 (moderate).
\begin{longtable}{P{0.15\linewidth}P{0.75\linewidth}}
\toprule Extraction field & Recorded observation\\\midrule\endhead
Mechanism & Temporal SLO-aware scheduling + feedback goodput balancing; MPS comparator\\
Environment & Dell PowerEdgeR740; V10032GB; local synthetic co-running experiment and separate four-GPU realistic trace. OS/driver/CUDA/backend versions NOT EXTRACTED.\\
Workload & Local GoogLeNet/ResNet152 Poisson arrivals; three-model test adds Mobile\_Pose\_MobileNetv3; 600req/s distributed in six ratios; SLO100ms and interval50ms. Separate61-model, 90min MAF replay on four GPUs.\\
Comparator & Clockwork; MPS plus Clockwork SLO scheduling with two worker processes/GPU; EDF without batching (not a like-for-like batching baseline).\\
Bounded finding & Local Table3 load3:1:2: unfairness1.05 vs Clockwork1.50; SLO goodput524.29 vs457.24req/s. Other load ratios give smaller differences. Local unfairness is max/min model slowdown(Eq2), not p99 or security. Four-GPU MAF outcome kept separate.\\
Transfer limit & Two-column Markdown interleaves algorithms/prose; use labelled Table3 rows and sections, not scrambled formulas. Repetitions/variance and software versions unresolved. MAF trace has minute-level invocation counts; cross-study batching and isolation definitions differ.\\
Source anchor & Published pp17-18; sections5.1-5.4; Table3; Eq2; Markdown lines718-899\\
\bottomrule\end{longtable}
\begin{longtable}{P{0.09\linewidth}P{0.10\linewidth}P{0.67\linewidth}}
\toprule QA & Score & Source-based rationale\\\midrule\endhead
Q1 & 1 & Heterogeneous model-sharing fairness explicit\\
Q2 & 1 & Estimator, feedback controller and temporal scheduling described\\
Q3 & 0.5 & Server/GPU/memory clear; OS/driver/backend unresolved\\
Q4 & 1 & Models, Poisson load, SLO, interval, load ratios and replay duration explicit\\
Q5 & 0.5 & Clockwork and adapted MPS; EDF lacks batching and is less directly comparable\\
Q6 & 1 & Goodput requests/s and unfairness ratio distinguish outcome definitions\\
Q7 & 0.5 & 90min replay and parameter sensitivity clear; local repetitions/variance unresolved\\
Q8 & 1 & Table3 rows/labels unambiguous in Markdown; plot coordinates absent\\
Q9 & 0.5 & Future work/memory and heterogeneity discussed; incomplete validity treatment\\
Q10 & 0.5 & Algorithms interpretable despite column interleaving; software/control gaps restrict replication\\
\bottomrule\end{longtable}

